% =============================================================
% Relativistic Viscous Rayleigh--Taylor Instability (Ch X, §94)
% Agent 44: Israel-Stewart causal viscosity extension
% =============================================================
%
% Classical source: output/chapters/chapter_10.tex, lines 577--947
% Conventions: RELATIVISTIC_CONVENTIONS.md, rel_preamble.tex
%
% This file extends Chandrasekhar's §94 analysis of two uniform
% viscous fluids separated by a horizontal boundary (Rayleigh--Taylor
% problem) to the relativistic regime, using Israel-Stewart causal
% dissipation for the viscous sector.
%
% Key physical points:
%   (i)    The Israel-Stewart relaxation time tau_pi renders viscous
%          signal propagation causal (finite speed), unlike the
%          acausal Navier-Stokes limit.
%   (ii)   The enthalpy density w = (epsilon + p)/c^2 replaces the
%          mass density rho in the inertial terms.
%   (iii)  Surface tension is generalised to a relativistic surface
%          energy density.
%   (iv)   Gravity waves acquire relativistic corrections through w.
%   (v)    All results reduce to Chandrasekhar's in the limit
%          c -> infinity, tau_pi -> 0.
% =============================================================

\section{Relativistic viscous Rayleigh--Taylor instability (\S94 -- relativistic)}
\label{sec:rel-10-94}

We now extend the classical analysis of \S\,94 --- two superposed
uniform viscous fluids separated by a horizontal boundary --- to the
special-relativistic regime.  The essential new ingredients are:
(a)~the replacement of the mass density $\rho$ by the relativistic
enthalpy density $\enthalpy = (\varepsilon + p)/c^{2}$ wherever
inertia enters; (b)~the use of the Israel--Stewart causal transport
equations for the viscous stress, which introduces a finite relaxation
time $\taupi > 0$; and (c)~the generalisation of surface tension $T$
to a relativistic surface energy density.

Throughout this section we work in the local rest frame of the
unperturbed interface and assume that the Lorentz factors $\Lf_{1,2}$
of the equilibrium fluids are unity (both fluids at rest).  This is
the natural relativistic generalisation of Chandrasekhar's set-up.
We keep $c$ explicit so that every non-relativistic limit is
transparent.

% -----------------------------------------------------------------
\subsection{Israel--Stewart viscous perturbation equations}
\label{sec:rel-10-94-IS}

In each uniform fluid region the linearised Israel--Stewart equation
for the shear-stress perturbation $\delta\pi^{\mu\nu}$ reads
(cf.\ the conventions in \S\,33 of \texttt{RELATIVISTIC\_CONVENTIONS.md})
\begin{equation}\label{eq:rel-10-IS-shear}
  \taupi\,\frac{\partial}{\partial t}\,\delta\pi^{\mu\nu}
  + \delta\pi^{\mu\nu}
  = 2\,\shearvisc\,\delta\sigma^{\mu\nu},
\end{equation}
where $\taupi > 0$ is the viscous relaxation time and $\shearvisc$ is
the shear viscosity.  For a perturbation $\propto e^{nt}$ (with $n$
the growth rate) this becomes
\begin{equation}\label{eq:rel-10-IS-shear-mode}
  \delta\pi^{\mu\nu}
  = \frac{2\,\shearvisc}{1 + \taupi\,n}\;\delta\sigma^{\mu\nu}
  \;\equiv\; 2\,\shearvisc^{\mathrm{eff}}(n)\;\delta\sigma^{\mu\nu}.
\end{equation}
The \emph{effective viscosity}
\begin{equation}\label{eq:rel-10-eta-eff}
  \shearvisc^{\mathrm{eff}}(n)
  = \frac{\shearvisc}{1 + \taupi\,n}
\end{equation}
is frequency-dependent; in the limit $\taupi \to 0$ we recover the
instantaneous (Navier--Stokes) value $\shearvisc$.

For the bulk viscous pressure we similarly write
\begin{equation}\label{eq:rel-10-IS-bulk}
  \tau_{\Pi}\,\frac{\partial}{\partial t}\,\delta\Pi
  + \delta\Pi
  = -\bulkvisc\;\nabla_{\mu}\,\delta u^{\mu},
\end{equation}
where $\tau_{\Pi} > 0$ is the bulk relaxation time and $\bulkvisc$
is the bulk viscosity coefficient.

% -----------------------------------------------------------------
\subsection{Relativistic dispersion equation in each uniform region}
\label{sec:rel-10-94-disp}

In the classical analysis the perturbation equation in each region of
constant $\rho$ and $\mu$ reduces to equation~\eqref{eq:10-96},
\begin{equation}\tag{\ref{eq:10-96}}
  \Bigl[1 - \frac{\nu}{n}\,(D^{2}-k^{2})\Bigr](D^{2}-k^{2})\,w = 0.
\end{equation}
In the relativistic theory the mass density $\rho$ is replaced by
the enthalpy density $\enthalpy = (\varepsilon+p)/c^{2}$, and the
kinematic viscosity is replaced by the \emph{effective} kinematic
viscosity
\begin{equation}\label{eq:rel-10-nu-eff}
  \nu^{\mathrm{eff}}(n)
  = \frac{\shearvisc^{\mathrm{eff}}(n)}{\enthalpy}
  = \frac{\nu}{1 + \taupi\,n},
  \qquad
  \nu \equiv \frac{\shearvisc}{\enthalpy}.
\end{equation}
The relativistic analogue of~\eqref{eq:10-96} is therefore
\begin{equation}\label{eq:rel-10-96}
  \boxed{
  \Bigl[1 - \frac{\nu}{n(1+\taupi\,n)}\,(D^{2}-k^{2})\Bigr]
  (D^{2}-k^{2})\,w = 0.
  }
\end{equation}
Writing $q_{\mathrm{rel}}^{2}$ for the non-trivial root of the
operator, we obtain the relativistic generalisation
of~\eqref{eq:10-98}:
\begin{equation}\label{eq:rel-10-98}
  q_{\mathrm{rel}}^{2}
  = k^{2} + \frac{n(1+\taupi\,n)}{\nu}
  = k^{2} + \frac{n}{\nu} + \frac{\taupi\,n^{2}}{\nu}.
\end{equation}
The solutions in the two half-spaces retain the same exponential
structure~\eqref{eq:10-99}--\eqref{eq:10-100} with $q_{1,2}$
replaced by $q_{\mathrm{rel},1,2}$.

\paragraph{Remark.}
In the Navier--Stokes limit $\taupi \to 0$, equation~\eqref{eq:rel-10-98}
reduces to $q^{2} = k^{2}+n/\nu$, recovering~\eqref{eq:10-98}.

% -----------------------------------------------------------------
\subsection{Boundary conditions at the relativistic interface}
\label{sec:rel-10-94-BC}

The boundary conditions at $z = 0$ --- continuity of $w$ and $Dw$,
continuity of tangential stress, and the normal-stress balance ---
carry over with two modifications:

\begin{enumerate}
  \item \emph{Inertial terms.} Every occurrence of $\rho_{1,2}$ in
    the inertial parts of the classical conditions
    \eqref{eq:10-104}--\eqref{eq:10-107} is replaced by
    $\enthalpy_{1,2} = (\varepsilon_{1,2}+p_{1,2})/c^{2}$.
  \item \emph{Viscous terms.} The dynamic viscosity
    $\mu_{i} = \rho_{i}\nu_{i}$ is replaced by
    $\enthalpy_{i}\,\nu_{i}^{\mathrm{eff}}(n)
    = \shearvisc_{i}/(1+\taupi_{i}\,n)$.
  \item \emph{Surface tension.} The classical surface tension $T$
    (force per unit length) is replaced by the relativistic surface
    energy density
    \begin{equation}\label{eq:rel-10-surface-energy}
      \mathcal{T} = T\Bigl(1 + \frac{T}{2\,\Sigma\,c^{2}}\Bigr),
    \end{equation}
    where $\Sigma$ is the surface mass density of the interface.
    For $c \to \infty$, $\mathcal{T} \to T$.
\end{enumerate}

The normal-stress balance (cf.\ the classical
equation~\eqref{eq:10-103}) becomes
\begin{align}\label{eq:rel-10-103}
  -k\,\enthalpy_{2}\,A_{2} - k\,\enthalpy_{1}\,A_{1}
  &= \frac{g\,k^{2}}{n^{2}}\,
    \Bigl[(\enthalpy_{2}-\enthalpy_{1})
    - \frac{k^{2}\mathcal{T}}{g}\Bigr]\,w_{0}
    \nonumber\\
  &\quad + \frac{k^{2}}{n}\,
    \bigl(\shearvisc_{1}^{\mathrm{eff}}-\shearvisc_{2}^{\mathrm{eff}}\bigr)
    \,(Dw)_{0}.
\end{align}

% -----------------------------------------------------------------
\subsection{The case $\nu_{1} = \nu_{2}$: relativistic characteristic equation}
\label{sec:rel-10-94a}

Following Chandrasekhar's simplification to equal kinematic
viscosities, $\nu_{1} = \nu_{2} = \nu$, and equal relaxation times
$\taupi_{1} = \taupi_{2} = \taupi$, we set
\begin{equation}\label{eq:rel-10-114}
  q_{\mathrm{rel}}^{2}
  = k^{2} + \frac{n(1+\taupi\,n)}{\nu},
\end{equation}
with the same $q_{\mathrm{rel}}$ in both half-spaces.  Define the
relativistic density fractions
\begin{equation}\label{eq:rel-10-108}
  \alpha_{1} = \frac{\enthalpy_{1}}{\enthalpy_{1}+\enthalpy_{2}},
  \qquad
  \alpha_{2} = \frac{\enthalpy_{2}}{\enthalpy_{1}+\enthalpy_{2}},
  \qquad
  \alpha_{1}+\alpha_{2} = 1,
\end{equation}
where $\enthalpy_{i} = (\varepsilon_{i}+p_{i})/c^{2}$.

The derivation of the characteristic equation parallels that of
\S\,94\,(a) exactly, with the substitutions
$\rho_{i} \to \enthalpy_{i}$ and $q \to q_{\mathrm{rel}}$.  The
result is the relativistic analogue of the quartic~\eqref{eq:10-121}.
Introducing
\begin{equation}\label{eq:rel-10-117}
  z_{\mathrm{rel}}
  = \frac{n(1+\taupi\,n)}{k^{2}\nu},
  \qquad
  y = \sqrt{1+z_{\mathrm{rel}}} = q_{\mathrm{rel}}/k,
\end{equation}
and the relativistic parameters
\begin{equation}\label{eq:rel-10-120}
  Q_{\mathrm{rel}}
  = \frac{g}{k^{3}\nu^{2}},
  \qquad
  S_{\mathrm{rel}}
  = \frac{\mathcal{T}}{(\enthalpy_{1}+\enthalpy_{2})\,\nu^{4/3}},
\end{equation}
the characteristic equation becomes
\begin{equation}\label{eq:rel-10-121}
  \boxed{
  \begin{aligned}
  &y^{4}+4\alpha_{1}\alpha_{2}\,y^{3}
    + 2(1-6\alpha_{1}\alpha_{2})\,y^{2}
    - 4(1-3\alpha_{1}\alpha_{2})\,y \\
  &\quad + (1-4\alpha_{1}\alpha_{2})
    + Q_{\mathrm{rel}}\,(\alpha_{1}-\alpha_{2})
    + Q_{\mathrm{rel}}^{4/3}\,S_{\mathrm{rel}} = 0.
  \end{aligned}
  }
\end{equation}
This is \emph{structurally identical} to Chandrasekhar's
equation~\eqref{eq:10-121}, but the variables carry different
physical content: $\alpha_{i}$ are enthalpy fractions (not mass
fractions), $S_{\mathrm{rel}}$ involves the relativistic surface
energy density, and -- crucially -- the relation between $y$ and the
growth rate $n$ is modified by the Israel--Stewart term:
\begin{equation}\label{eq:rel-10-n-from-y}
  y^{2} - 1
  = \frac{n(1+\taupi\,n)}{k^{2}\nu}
  \quad\Longrightarrow\quad
  \taupi\,n^{2} + n - k^{2}\nu\,(y^{2}-1) = 0.
\end{equation}
Solving for $n$:
\begin{equation}\label{eq:rel-10-n-quadratic}
  n = \frac{-1 + \sqrt{1 + 4\,\taupi\,k^{2}\nu\,(y^{2}-1)}}
           {2\,\taupi}.
\end{equation}
In the limit $\taupi \to 0$ this reduces to
$n = k^{2}\nu\,(y^{2}-1)$, recovering~\eqref{eq:10-124}.

\paragraph{Physical interpretation.}
For a given root $y$ of~\eqref{eq:rel-10-121}, the growth rate $n$
is \emph{reduced} relative to the Navier--Stokes prediction
(at fixed $y$) by the factor $(1+\taupi\,n)^{-1}$ implicit
in~\eqref{eq:rel-10-n-quadratic}.  The Israel--Stewart relaxation
effectively ``stiffens'' the viscous response at high frequencies,
slowing the exponential growth of short-wavelength modes.

% -----------------------------------------------------------------
\subsection{Modes of maximum instability: relativistic corrections}
\label{sec:rel-10-94b}

When $\mathcal{T} = 0$ (no surface tension) and
$\alpha_{2} > \alpha_{1}$ (heavier fluid on top), the quartic
\eqref{eq:rel-10-121} again admits a single root with positive real
part, and the arrangement is unstable for all wavenumbers.

The asymptotic behaviour parallels the classical case:
\begin{itemize}
  \item \emph{Long wavelengths} ($k \to 0$, $y \to \infty$):
    $n \sim [(\alpha_{2}-\alpha_{1})\,g\,k]^{1/2}$ as in the
    inviscid (and non-relativistic) limit, since viscosity and
    Israel--Stewart corrections become negligible.  The only
    relativistic effect is through $\alpha_{i}$, which are enthalpy
    fractions rather than mass fractions.
  \item \emph{Short wavelengths} ($k \to \infty$, $y \to 1$): from
    \eqref{eq:rel-10-n-quadratic} with $y^{2}-1 \to 0$,
    \begin{equation}\label{eq:rel-10-131}
      n \;\to\; \frac{(\alpha_{2}-\alpha_{1})\,g}{2\,k\,\nu}
      \cdot\frac{1}{1 + \taupi\,(\alpha_{2}-\alpha_{1})\,g/(2\,k\,\nu)}.
    \end{equation}
    For $\taupi = 0$ this reduces to~\eqref{eq:10-131}.  For
    $\taupi > 0$ the growth rate is further suppressed at large $k$,
    confirming that the Israel--Stewart relaxation damps the
    shortest-wavelength modes.
\end{itemize}

Since $n \to 0$ both as $k \to 0$ and $k \to \infty$, a mode of
\emph{maximum instability} persists.  Its wavenumber $k_{\max}$ and
growth rate $n_{\max}$ are shifted relative to the classical values
listed in Table~\ref{tab:XLVI}: specifically, $k_{\max}$ moves to
\emph{smaller} $k$ (longer wavelength) and $n_{\max}$ is
\emph{reduced}, with corrections of relative order
$\taupi\,(g^{2}/\nu)^{1/3}$.  In the non-relativistic limit
$c \to \infty$ (whence $\enthalpy_{i} \to \rho_{i}$ and
$\taupi \to 0$), the classical Table~\ref{tab:XLVI} is exactly
recovered.

% -----------------------------------------------------------------
\subsection{Surface tension: relativistic surface energy density}
\label{sec:rel-10-94c}

When the relativistic surface energy density
$\mathcal{T}$~\eqref{eq:rel-10-surface-energy} is included, the
critical wavenumber above which surface tension stabilises the
instability is
\begin{equation}\label{eq:rel-10-kc}
  k_{c}
  = \sqrt{\frac{(\enthalpy_{2}-\enthalpy_{1})\,g}{\mathcal{T}}}
  = \sqrt{\frac{(\enthalpy_{2}-\enthalpy_{1})\,g}{T}}
    \;\Bigl(1 - \frac{T}{4\,\Sigma\,c^{2}}
    + \mathcal{O}(c^{-4})\Bigr).
\end{equation}
This reduces to the classical value~\eqref{eq:10-135} (with
$\rho_{i} \to \enthalpy_{i}$) in the limit $c \to \infty$.

The relativistic correction lowers $k_{c}$ slightly: the surface
energy contribution to the interface rest mass increases the effective
inertia, allowing marginally longer wavelengths to be stabilised.  As
in the classical case, $k_{c}$ is \emph{independent of viscosity} and
of the Israel--Stewart relaxation time $\taupi$.

For the $(n,k)$-relationships when $\mathcal{T} \neq 0$, the
analysis proceeds exactly as in \S\,94\,(c) with
$Q_{\mathrm{rel}}^{4/3}\,S_{\mathrm{rel}}$ replacing
$Q^{4/3}S$ in equation~\eqref{eq:10-133} and the growth rate
recovered from~\eqref{eq:rel-10-n-quadratic} rather than
from~\eqref{eq:10-124}.

% -----------------------------------------------------------------
\subsection{Gravity waves: relativistic dispersion relation}
\label{sec:rel-10-94e}

Passing to the limit $\alpha_{2} = 0$, $\alpha_{1} = 1$ (a single
fluid with a free surface, i.e.\ the ``deep ocean'' limit), the
relativistic characteristic equation~\eqref{eq:rel-10-121} becomes
\begin{equation}\label{eq:rel-10-136}
  y^{4} + 2\,y^{2} - 4\,y + 1
  + Q_{\mathrm{rel}} + Q_{\mathrm{rel}}^{4/3}\,S_{\mathrm{rel}} = 0.
\end{equation}
This is structurally Chandrasekhar's equation~\eqref{eq:10-136} but
with $Q_{\mathrm{rel}} = g/(k^{3}\nu^{2})$ and the enthalpy-based
$S_{\mathrm{rel}}$.

In the \emph{inviscid, zero-surface-tension} limit ($\nu \to 0$,
$\mathcal{T} = 0$), the long-wavelength gravity-wave dispersion
relation becomes
\begin{equation}\label{eq:rel-10-grav-disp}
  \omega^{2} = g\,k\,\frac{\rho}{\enthalpy}
  = g\,k\,\frac{\rho\,c^{2}}{\varepsilon + p}
  = g\,k\;\Bigl(1 - \frac{p + \varepsilon_{\mathrm{int}}}{\varepsilon+p}\Bigr),
\end{equation}
where $\varepsilon_{\mathrm{int}} = \varepsilon - \rho c^{2}$ is the
internal energy density.  For a non-relativistic fluid
$p \ll \rho c^{2}$ and $\varepsilon_{\mathrm{int}} \ll \rho c^{2}$,
so $\omega^{2} \to g\,k$, recovering the classical deep-water result.
In the ultra-relativistic limit ($p \to \varepsilon/3$, radiation
fluid), $\omega^{2} = (3/4)\,g\,k$: the increased inertia of the
relativistic enthalpy reduces the wave frequency.

When viscosity is retained, the branch-point structure (where damping
changes from aperiodic to oscillatory, cf.\ Table~\ref{tab:XLIX}) is
modified:
\begin{equation}\label{eq:rel-10-branch-shift}
  k_{a}^{\mathrm{rel}}
  = k_{a}^{\mathrm{class}}
    \Bigl[1 + \mathcal{O}\!\bigl(\taupi\,(g^{2}/\nu)^{1/3}\bigr)
    \Bigr],
\end{equation}
with the correction pushing the branch point to larger $k$ (shorter
wavelengths must be reached before the overdamped regime sets in).

% -----------------------------------------------------------------
\subsection{Causality of the Israel--Stewart viscous modes}
\label{sec:rel-10-94-causality}

In the classical (Navier--Stokes) treatment the viscous mode
$e^{qz}$ with $q^{2} = k^{2}+n/\nu$ has no upper bound on the
speed of viscous signal propagation: the parabolic diffusion equation
admits instantaneous propagation, violating relativistic causality.

The Israel--Stewart theory cures this defect.  The effective
dispersion relation~\eqref{eq:rel-10-98},
\[
  q_{\mathrm{rel}}^{2}
  = k^{2} + \frac{n}{\nu} + \frac{\taupi\,n^{2}}{\nu},
\]
is \emph{hyperbolic} in the high-frequency limit $|n| \to \infty$:
\begin{equation}\label{eq:rel-10-vvisc}
  q_{\mathrm{rel}}^{2}
  \;\xrightarrow{|n|\to\infty}\;
  \frac{\taupi}{\nu}\,n^{2} + k^{2},
\end{equation}
which describes propagation at the \emph{finite} viscous signal speed
\begin{equation}\label{eq:rel-10-vvisc-speed}
  \boxed{
  v_{\mathrm{visc}}
  = \sqrt{\frac{\nu}{\taupi}}
  = \sqrt{\frac{\shearvisc}{\enthalpy\,\taupi}}.
  }
\end{equation}
Causality requires $v_{\mathrm{visc}} \leq c$, which imposes the
constraint
\begin{equation}\label{eq:rel-10-causality-bound}
  \taupi \;\geq\; \frac{\shearvisc}{\enthalpy\,c^{2}}
  = \frac{\nu}{c^{2}}.
\end{equation}
This is precisely the causality bound derived by Israel and Stewart
(1979).  So long as~\eqref{eq:rel-10-causality-bound} is satisfied,
the viscous perturbation modes propagate at subluminal speeds and the
initial-value problem for the Rayleigh--Taylor instability is
well-posed in the relativistic sense.

In the non-relativistic limit $c \to \infty$ the
bound~\eqref{eq:rel-10-causality-bound} is trivially satisfied for
any $\taupi > 0$, and $v_{\mathrm{visc}} \to \infty$ as
$\taupi \to 0$, recovering the instantaneous diffusion of the
classical theory.

\paragraph{Summary.}
The Israel--Stewart formulation achieves three things simultaneously:
(i)~it removes the acausal instantaneous propagation of the
Navier--Stokes viscous modes; (ii)~it provides a natural
ultraviolet regularisation of the Rayleigh--Taylor instability at
short wavelengths via the frequency-dependent effective
viscosity~\eqref{eq:rel-10-eta-eff}; and (iii)~it preserves the
full structure of Chandrasekhar's classical results in the limit
$c \to \infty$, $\taupi \to 0$.

% -----------------------------------------------------------------
\subsection{Non-relativistic limit}
\label{sec:rel-10-94-NR}

We verify explicitly that all relativistic expressions reduce to
Chandrasekhar's classical results.  Setting
$c \to \infty$:
\begin{enumerate}
  \item $\enthalpy_{i} = (\varepsilon_{i}+p_{i})/c^{2}
    \to \rho_{i}$ (rest-mass density dominates).
  \item $\alpha_{i} \to \rho_{i}/(\rho_{1}+\rho_{2})$
    (mass fractions).
  \item $\taupi \to 0$ (relaxation time vanishes; Navier--Stokes
    recovered).
  \item $\nu^{\mathrm{eff}}(n) \to \nu$
    (frequency-independent viscosity).
  \item $q_{\mathrm{rel}} \to q = \sqrt{k^{2}+n/\nu}$
    (classical~\eqref{eq:10-98}).
  \item $\mathcal{T} \to T$ (classical surface tension).
  \item Equation~\eqref{eq:rel-10-121} $\to$
    equation~\eqref{eq:10-121} (classical quartic).
  \item $n = k^{2}\nu(y^{2}-1)$ (classical~\eqref{eq:10-124}).
\end{enumerate}
Every table, figure, and asymptotic formula of \S\,94 is thus
faithfully recovered.
